% LỜI TỰA --- Quyển XII
\chapter*{Lời Tựa}
\addcontentsline{toc}{chapter}{Lời Tựa}
\markboth{Lời Tựa}{Lời Tựa}

\begin{truyen}{Lời Tựa}{What the Teacher Sees}
\chuhoa{G}{iáo viên bắt đầu bộ sách này với một câu hỏi đơn giản: làm thế nào để mỗi buổi học có thể để lại thứ gì đó đáng nhớ hơn là bài thi?}
\textit{(A teacher began this series with a simple question: how can each class leave behind something more memorable than a test?)}

Bục giảng nhìn thấy những gì học sinh chưa tự nhìn thấy ở chính mình.
\textit{(The chalkboard sees what students have not yet seen in themselves.)}

Quyển mười hai tập hợp những câu chuyện từ góc nhìn của người đứng trên bục giảng~--- không phải để khoe sự quan sát, mà để chia sẻ những gì thầy giáo thực sự muốn học trò của mình hiểu. Đây là quyển sách thầy muốn viết cho trò. Và trò nên đọc với lòng biết ơn.
\textit{(Volume twelve collects stories from the perspective of those at the chalkboard~--- not to display their observations, but to share what teachers truly want their students to understand. This is the book a teacher wants to write for students. And students should read it with gratitude.)}

Bộ sách <<Truyện Tích Cảm Hứng Song Ngữ>> gồm 24 quyển, mỗi quyển một chủ đề riêng. Quyển XII này mang chủ đề: \textbf{góc nhìn từ bục giảng về học sinh}.
\textit{(The series ``Inspiring Bilingual Tales'' contains 24 volumes, each with its own theme. Volume XII focuses on: \textbf{góc nhìn từ bục giảng về học sinh}.})
\end{truyen}

\begin{baihoc}
Thầy giáo không phải là người chỉ dạy chữ~--- họ là người nhìn thấy tiềm năng trước khi học trò tự nhìn thấy.
\textit{(Teachers are not just those who teach letters~--- they are those who see potential before students see it themselves.)}
\end{baihoc}

\begin{ghinhoanh}
\textbf{Short English to remember:}

A teacher sees the student the student cannot see yet.

The classroom reveals what the world will later confirm.
\end{ghinhoanh}

\begin{tuhocnhanh}
1. Điều gì trong quyển này giúp em hiểu thầy cô hơn? \textit{(What in this volume helped you understand your teachers better?)}

2. Em từng được thầy cô nhìn thấy tiềm năng mình không? Lúc đó em cảm thấy thế nào? \textit{(Has a teacher ever seen potential in you that you didn't see? How did that feel?)}

3. Nếu em là thầy giáo, điều đầu tiên em muốn học trò mình hiểu là gì? \textit{(If you were a teacher, what is the first thing you would want students to understand?)}
\end{tuhocnhanh}

\begin{center}
\textit{\color{chalkgold}``Thầy nhìn thấy em~--- kể cả khi em chưa thấy chính mình.''}
\end{center}

\begin{flushright}
\textbf{Tôi Nguyễn Văn Sang}\\
\textit{GV THPT Nguyễn Hữu Cảnh - TP HCM}
\end{flushright}

\ngancach
